Indonesia menghasilkan sekitar 23-48 juta ton sampah makanan per tahun, dengan proporsi terbesar berasal dari sektor UMKM kuliner yang kehilangan 10-30\% produk akibat surplus yang tidak terjual. Masalah ini tidak hanya menimbulkan kerugian ekonomi bagi pelaku UMKM, tetapi juga berkontribusi signifikan terhadap emisi gas rumah kaca dan pemborosan sumber daya alam. Sayangnya, belum ada solusi terpadu yang secara efektif menjembatani UMKM kuliner dengan konsumen dan mitra donasi untuk menyalurkan makanan surplus yang masih layak konsumsi.

ResQfood hadir sebagai platform manajemen \textit{food waste} berbasis web yang menghubungkan UMKM kuliner dengan konsumen dan mitra donasi. Platform ini menawarkan serangkaian fitur inovatif meliputi \textit{Food Eligibility Assessment} untuk menilai kelayakan makanan secara objektif, \textit{Dynamic Pricing} untuk mengoptimalkan harga jual berdasarkan kondisi makanan, \textit{Sistem Prediksi Produksi Harian} berbasis data historis dan cuaca, serta mekanisme donasi otomatis melalui \textit{Food Rescue Score}. Dengan memanfaatkan teknologi \textit{machine learning} untuk prediksi produksi dan analitik \textit{real-time}, ResQfood bertujuan mengurangi \textit{food waste} hingga 40\% sekaligus meningkatkan pendapatan UMKM melalui penjualan makanan surplus.

Platform ini memberikan dampak positif bagi tiga pihak utama: UMKM mendapatkan pendapatan tambahan dari makanan yang sebelumnya terbuang, konsumen memperoleh akses terhadap makanan berkualitas dengan harga terjangkau, serta lingkungan mendapat manfaat dari pengurangan emisi gas rumah kaca dan volume sampah organik. Dengan target implementasi di Kota Bandung pada tahun pertama, ResQfood berpotensi menjadi solusi \textit{scalable} yang dapat diterapkan di kota-kota besar lainnya di Indonesia.
